% !TeX root = ../PIFB2.tex
\section{Exact Coarse-Graining and Induced Interactions}\label{sec:rg_status}

Coarse-graining can generate cross-scale operators without postulating their form. It does so once a fine law and a channel have been supplied. This conditional statement is stronger than a formal analogy with renormalization and weaker than a derivation of a preferred hierarchy.

\subsection{The effective action}
Keep the coarse channel $C$ fixed when varying the recognition law. Let $\rho(dy)$ be a proper reference law and $m_o(dy)=e^{-H_o(y)}\rho(dy)$ a finite positive evidence measure. Push both through the same channel:
\begin{equation}
\rho^c=\rho C,\qquad m_o^c=m_oC,\qquad
H_o^c(z)=-\log\frac{dm_o^c}{d\rho^c}(z).
\label{eq:effective_action}
\end{equation}
A reverse conditional under $\rho(dy)C(y,dz)$ gives
\begin{equation}
e^{-H_o^c(z)}=\mathbb E_\rho[e^{-H_o(Y)}\mid Z=z].
\label{eq:rg_conditional_partition}
\end{equation}
To prove it, integrate both sides against an arbitrary bounded measurable function of $z$. Both integrals equal $\int f(z)e^{-H_o(y)}\rho(dy)C(y,dz)$. This identifies the Radon--Nikodym derivative. Total mass, hence the evidence, is preserved. Successive channels compose exactly when the conditionals are taken from one joint tower~\cite[Ch.~11]{DennisGaugeVFE2026}.

When $\rho$ is a latent prior and $e^{-H_o}$ a record likelihood, this is an evidence-preserving coarse model. When $\rho$ is instead a reference on belief configurations and $H=\mathcal E/T_{\mathrm{cfg}}$, it is a different probability model over those configurations. Deterministic descent of $\mathcal E$ does not sample that Gibbs ensemble. The configuration temperature $T_{\mathrm{cfg}}$, the source temperature $\tau$ and physical thermodynamic temperature are distinct.

\subsection{What integration generates}
For a bounded perturbation $V$ of $H_o$, restrict to coarse states where $0<dm_o^c/d\rho^c<\infty$, which holds $P_o^c$-almost surely. The effective perturbation is
\begin{equation}
H_{o,\lambda}^c-H_o^c
=-\log\mathbb E_{P_o(dy\mid z)}e^{-\lambda V(Y)}
=\lambda\mathbb E[V\mid z]
-\frac{\lambda^2}{2}\operatorname{Var}(V\mid z)+O(\lambda^3).
\label{eq:conditional_cumulant_action}
\end{equation}
The expansion follows by differentiating the conditional moment generating function near zero; boundedness supplies domination and local analyticity. For sums of fine interactions the variance contains cross covariances, and higher cumulants couple additional retained variables. Thus exact coarse-graining generally generates higher-order interactions. Assuming the same pairwise functional at every scale discards these terms unless a closure theorem applies.

An elementary Gaussian calculation makes the induced operators explicit. Suppose the dimensionless quadratic energy is
\begin{equation}
H(x,z)=\frac12x^\top Ax+x^\top Bz+\frac12z^\top Dz,\qquad A\succ0.
\label{eq:quadratic_elimination}
\end{equation}
Integration over $x$ with Lebesgue reference gives
\begin{equation}
H^c(z)=\frac12z^\top(D-B^\top A^{-1}B)z
+\frac12\log\det A-\frac{d_x}{2}\log(2\pi).
\label{eq:schur_effective_action}
\end{equation}
Completing the square proves the formula. For a proper joint Gaussian, the Schur complement is positive definite. If $A=A(z)$, its log determinant is an induced potential on retained states and must remain in the action. Completing the square still works pointwise when the integration is finite, but the result need not be quadratic.

The matrix $B^\top A^{-1}B$ supplies interactions between retained variables that previously communicated only through the eliminated block. It can also generate self terms. A closure family that forbids every newly generated self term is narrower than the family of valid Gaussian effective laws. The companion Gaussian theory makes this distinction explicit~\cite{DennisGaugeVFESuppG2026}. These induced terms are examples of cross-scale structure obtained by integration, not additional forces appended to the fine functional.

\subsection{Closure and its residual}
If the exact coarse action $H^c$ is replaced by an ansatz $\widehat H^c$, write $R=H^c-\widehat H^c$. Against the same reference, the exact law satisfies
\begin{equation}
p^c=\frac{e^{-R}\widehat p^c}{\mathbb E_{\widehat p^c}e^{-R}}.
\label{eq:coarse_residual_law}
\end{equation}
For a bounded residual with $\sup R-\inf R\leq\delta$,
\begin{equation}
D_{\mathrm{KL}}(p^c\|\widehat p^c)
=-\mathbb E_{p^c}R-\log\mathbb E_{\widehat p^c}e^{-R}
\leq\delta.
\label{eq:closure_residual_bound}
\end{equation}
The inequality follows by bounding the logarithmic normalizer between $-\sup R$ and $-\inf R$. Additive constants in the action cancel. This elementary bound is conservative, but it measures a stated approximation; a small visual mismatch of Gaussian means is not an equivalent criterion.

A finite sequence of channels is exact coarse-graining. To obtain a renormalization flow on one coupling space, one must additionally identify retained spaces, specify rescaling and control the generated operator family. Finiteness alone does not prohibit defining comparison maps or fixed points; it means that they are not supplied by decreasing the number of agents. The finite-step theory and the spatial or infinite-system extensions answer different questions. The latter require their additional hypotheses~\cite{DennisGaugeVFESuppB2026}.

\subsection{Persistent laws versus instantaneous marginals}
An exact marginal need not have a closed Markov evolution. If the fine generator is $\mathcal L$ and retention is deterministic, a sufficient generator closure is
\begin{equation}
\mathcal L(f\circ C)=(\mathcal L^cf)\circ C
\label{eq:generator_descent}
\end{equation}
on a determining domain, together with the well-posedness assumptions that turn generators into the stated processes. If fine states sharing $C(y)$ predict different future retained laws, an instantaneous parent state is insufficient.

The conditional law of the hidden fine state given retained history is the canonical exact repair. It is generally a measure-valued state; finite-dimensional closure is a property to prove for the particular system. A generic dimension-counting argument cannot exclude arbitrary measurable encodings, and ordinary path relative entropy may be infinite when candidate processes have different quadratic variation. Approximation must therefore name its class, observations and time horizon. Section~\ref{sec:timescales} and Appendix~\ref{app:examples} develop constructive memory calculations rather than assuming memory away.

